\documentclass[a4paper]{article}

% Basic packages
% Español (México) con XeLaTeX: fontspec para fuentes Unicode, polyglossia
% para la división silábica y los nombres del documento en la variante mexicana.
\usepackage{fontspec}
\usepackage{polyglossia}
\setdefaultlanguage[variant=mexican]{spanish}
% Los nombres chinos van como «pinyin (original)»: xeCJK da a los caracteres
% Han su propia fuente; sin ella XeLaTeX los omite con un aviso.
% FandolSong no cubre algunos caracteres de nombres (U+73FA); WenQuanYi Zen Hei sí.
\usepackage[AutoFallBack=true]{xeCJK}
\setCJKmainfont{FandolSong-Regular.otf}
\setCJKfallbackfamilyfont{\CJKrmdefault}{WenQuanYi Zen Hei}
% polyglossia no traduce los nombres de listings.
\AtBeginDocument{\providecommand{\lstlistingname}{}\renewcommand{\lstlistingname}{Listado}%
  \providecommand{\lstlistlistingname}{}\renewcommand{\lstlistlistingname}{Índice de listados}}
\usepackage{amsmath, amssymb}
\usepackage{graphicx}
\usepackage[margin=2.5cm]{geometry}
\usepackage[most]{tcolorbox}
\usepackage{etoolbox}
\usepackage{listings}
\usepackage{hyperref}
\usepackage{booktabs}
\usepackage{subcaption}
\usepackage{float}

% Highlight boxes
\newtcolorbox{knowledgebox}[1]{
    enhanced, colback=blue!5!white, colframe=blue!75!black, colbacktitle=blue!75!black,
    coltitle=white, fonttitle=\bfseries, title={#1}, attach boxed title to top left={yshift=-2mm, xshift=2mm},
    boxrule=1pt, sharp corners
}

\newtcolorbox{importantbox}[1]{
    enhanced, colback=yellow!10!white, colframe=yellow!80!black, colbacktitle=yellow!80!black,
    coltitle=black, fonttitle=\bfseries, title={#1}, sharp corners
}

\newtcolorbox{warningbox}[1]{
    enhanced, colback=red!5!white, colframe=red!75!black, colbacktitle=red!75!black,
    coltitle=white, fonttitle=\bfseries, title={#1}, sharp corners
}

% Listings
\lstset{
    language=Python,
    basicstyle=\ttfamily\small,
    keywordstyle=\color{blue},
    stringstyle=\color{red!60!black},
    commentstyle=\color{green!60!black},
    breaklines=true,
    frame=single,
    numbers=left,
    numberstyle=\tiny\color{gray},
    captionpos=b,
    extendedchars=true
}

% Front-page metadata
\newcommand{\notetitle}{Entrevista a Yang Zhilin (杨植麟): el modelo K2, Scaling Law y el futuro del Agent}
\newcommand{\noteauthors}{Elaboradas a partir de materiales públicos del curso}
\newcommand{\notedate}{julio de 2025}
\newcommand{\videochannel}{Zhang Xiaojun Shangye Fangtanlu (张小珺商业访谈录)}
\newcommand{\videopublishdate}{julio de 2025}
\newcommand{\videoduration}{aproximadamente 90 minutos}
\newcommand{\videourl}{}
\newcommand{\videocoverpath}{cover.jpg}
\newcommand{\repourl}{https://github.com/hqhq1025/ai-course-notes}


% Footer with repo info
\usepackage{fancyhdr}
\pagestyle{fancy}
\fancyhf{}
\fancyfoot[C]{\thepage}
\fancyfoot[R]{\footnotesize\color{gray}\href{\repourl}{AI Course Notes}}
\renewcommand{\headrulewidth}{0pt}
\renewcommand{\footrulewidth}{0pt}

\begin{document}

\begin{titlepage}
\centering
{\Large Notas de entrevista\par}
\vspace{1.2cm}
{\huge\bfseries \notetitle\par}
\vspace{0.8cm}
{\large \noteauthors\par}
\vspace{0.3cm}
{\large \notedate\par}
\vspace{1.2cm}

\ifdefempty{\videocoverpath}{
{\small Al generar las notas, indique la ruta local de la portada del video.\par}
}{
\includegraphics[width=0.82\textwidth,height=0.45\textheight,keepaspectratio]{\videocoverpath}\par
}

\vfill
\begin{tcolorbox}[width=0.9\textwidth, colback=black!2!white, colframe=black!60, sharp corners]
\textbf{Autor o canal del video}: \videochannel\par
\textbf{Fecha de publicación}: \videopublishdate\par
\textbf{Duración del video}: \videoduration\par
\ifdefempty{\videourl}{}{
\textbf{Enlace del video}: \href{\videourl}{\nolinkurl{\videourl}}\par
}
\par
\textbf{Página del proyecto}: \href{\repourl}{\nolinkurl{\repourl}}\par
\end{tcolorbox}
\end{titlepage}

\tableofcontents
\newpage

%% --- Inicio del contenido --- %%

\section{Escalar la montaña nevada sin fin: el camino interior de emprender en IA}

\subsection{De «avanzar hacia una montaña nevada interminable y desconocida» a «Beginning of Infinity»}

Yang Zhilin (杨植麟) es fundador y director ejecutivo de Moonshot AI (月之暗面). Un año después (julio de 2025), en una nueva entrevista, describió lo que sentía recurriendo a la filosofía de \textit{The Beginning of Infinity}:

\begin{importantbox}{Los problemas son inevitables, pero se pueden resolver}
David Deutsch propone en \textit{The Beginning of Infinity} dos tesis centrales:
\begin{enumerate}
    \item Los problemas son inevitables (Problems are inevitable)
    \item Los problemas se pueden resolver (Problems are soluble)
\end{enumerate}
El proceso de investigación y desarrollo de IA encarna precisamente esta filosofía: se resuelven muchos problemas del aprendizaje por refuerzo y surgen problemas nuevos de evaluación y verificación; se resuelve el problema de la eficiencia del preentrenamiento y surge el problema de la generalización de los agentes. Cada vez que se resuelve un problema, la tecnología avanza unos cientos de metros más y, al mismo tiempo, aparecen problemas nuevos.
\end{importantbox}

Yang Zhilin considera que la AGI tal vez no sea un escalón determinado, sino una dirección. La IA actual ya hace mejor que el 99\% de los seres humanos en muchos campos; en cierto sentido puede considerarse una AGI parcial. Pero es difícil decir que en un momento dado se alcanza de pronto la AGI: se parece más a un proceso de ascenso continuo.

\begin{knowledgebox}{Puede que la montaña nevada no tenga fin}
Yang Zhilin expresó una postura optimista: espera que esta montaña nevada nunca tenga fin: «eso es justo lo que significa Beginning of Infinity, puede que sea una montaña infinita». Y en cierta etapa, tal vez ya no sean los seres humanos quienes escalen, sino que se use la IA para escalar; por ejemplo, usar el modelo K2 para participar en el entrenamiento de modelos y en el procesamiento de datos. La IA se está convirtiendo en un «amplificador de la civilización humana».
\end{knowledgebox}

\subsection{Resumen de la sección}

Emprender en IA es un proceso continuo de resolver problemas y descubrir problemas nuevos. La AGI no es un punto final, sino una dirección. La IA está pasando de ser una herramienta a ser un participante en el proceso de investigación y desarrollo.

\section{Los cambios de paradigma en los grandes modelos globales durante el último año}

\subsection{Primer paradigma: modelos de razonamiento con pensamiento intensivo}

Los modelos de razonamiento (Reasoning Model) basados en aprendizaje por refuerzo, representados por O1, son uno de los avances técnicos más importantes del último año.

\begin{importantbox}{La esencia de los modelos de razonamiento: el ciclo de conjetura y verificación}
La capacidad central de los modelos de razonamiento es la «reflexión» (Reflection), que en esencia comprende dos capacidades:
\begin{enumerate}
    \item \textbf{Formular conjeturas}: proponer continuamente nuevas hipótesis durante el proceso de resolución del problema
    \item \textbf{Autoverificación}: verificar de manera implícita cada conjetura para determinar si es correcta o incorrecta
\end{enumerate}
Mediante múltiples ciclos de «conjetura → verificación → nueva conjetura → nueva verificación», el modelo puede convertir un problema de Pass@2 en uno de Pass@1: en esencia, es equivalente a haberlo intentado varias veces. Esto se parece mucho al proceso humano de investigación científica y resolución de problemas.
\end{importantbox}

Yang Zhilin (杨植麟) compara de manera gráfica este tipo de modelo de razonamiento con un \textbf{«cerebro en una cubeta»}: no necesita interactuar con el mundo exterior, sino que solo piensa dentro de su propio cerebro, y así puede resolver problemas mediante el ciclo de conjetura y verificación.

\begin{knowledgebox}{Muestreo serial vs. paralelo}
La forma efectiva de trabajo de los modelos de razonamiento es principalmente serial: cada conjetura se basa en el resultado de la anterior. Aunque también se puede combinar con una estrategia paralela (muestrear varias soluciones al mismo tiempo), las investigaciones recientes muestran que el límite superior de la estrategia serial suele ser más alto. Esto coincide con las conclusiones experimentales de Moonshot AI (月之暗面).
\end{knowledgebox}

\subsection{Segundo paradigma: el aprendizaje por refuerzo de Agent}

A diferencia del «cerebro en una cubeta» de los modelos de razonamiento, el modelo de Agent interactúa con el entorno exterior mediante múltiples turnos:

\begin{itemize}
    \item Piensa y ejecuta acciones al mismo tiempo (búsqueda, navegador, escritura de código, etc.)
    \item Resuelve problemas mediante un método de múltiples turnos
    \item La siguiente acción depende de la retroalimentación externa y de la actualización del estado
\end{itemize}

\begin{importantbox}{La unificación de dos formas de Test-Time Scaling}
Los modelos de razonamiento y los modelos de Agent apuntan al mismo concepto central: \textbf{Test-Time Scaling}:
\begin{itemize}
    \item \textbf{Modelos de razonamiento}: hacen Scale de la cantidad de tokens de pensamiento dentro de un solo turno (más razonamiento interno)
    \item \textbf{Modelos de Agent}: hacen Scale del número de turnos de interacción (más interacción externa)
\end{itemize}
Ambos consisten en invertir más recursos de cómputo en el momento de la inferencia para completar tareas más complejas. Un modelo de conversación tradicional puede generar solo unos cientos de tokens, mientras que un modelo de razonamiento o de Agent puede tardar varias horas y consumir una gran cantidad de tokens para completar una tarea compleja.
\end{importantbox}

\subsection{Claude vs. OpenAI: diferencias de ruta técnica}

\begin{knowledgebox}{Reasoning y Agent no tienen necesariamente una dependencia lineal}
Yang Zhilin señala una observación interesante: los modelos de Claude no destacan en los benchmarks de Reasoning, pero tienen un desempeño muy bueno en las tareas de Agent. Esto indica que:
\begin{itemize}
    \item Reasoning y Agent corresponden a dos dimensiones distintas de Test-Time Scaling
    \item En el desarrollo, no existe un orden secuencial necesario entre ambos
    \item Pero para alcanzar el nivel más alto de Agent, al final sigue siendo necesaria una capacidad de Reasoning muy sólida
\end{itemize}
Es posible que Claude haya invertido más en la dirección de Agent (interacción de múltiples turnos), mientras que OpenAI ha invertido más en la dirección de Reasoning (pensamiento profundo). Las distintas rutas técnicas generan diferencias de producto en el corto plazo, pero a largo plazo ambas son indispensables.
\end{knowledgebox}

\subsection{Tercera tendencia: las empresas de modelos hacen producto propio}

\begin{importantbox}{Diseño directo vs. ingeniería inversa}
Yang Zhilin distingue dos formas de desarrollar productos:
\begin{itemize}
    \item \textbf{Ingeniería inversa de terceros}: a partir de un modelo base existente, se hace ingeniería inversa del proceso de entrenamiento del modelo mediante el diseño de andamiajes, herramientas y System Prompt, para encontrar la mejor forma de usarlo (como Manus, entre otros)
    \item \textbf{Diseño directo propio}: la empresa de modelos primero diseña las herramientas y el método de Context Engineering, y después entrena el modelo en ese entorno, de modo que el modelo rinda mejor de forma natural en su propio entorno (como Claude Code, ChatGPT Agent)
\end{itemize}
La segunda forma probablemente tiene un límite superior más alto, porque permite integrar mejor las herramientas y el modelo, y entrenar de extremo a extremo.
\end{importantbox}

\subsection{Resumen de la sección}

Los tres grandes cambios de paradigma del último año son: (1) los modelos de razonamiento (el ciclo de conjetura y verificación al estilo del cerebro en una cubeta), (2) los modelos de Agent (múltiples turnos de interacción con el mundo exterior) y (3) las empresas de modelos que hacen producto propio. El razonamiento y Agent son dos dimensiones de Test-Time Scaling que, al final, deben combinarse.

\section{Una reconsideración de la clasificación L1-L5 de OpenAI}

\subsection{El significado de los cinco niveles}

OpenAI definió cinco niveles de capacidad de la IA: L1 ChatBot → L2 Reasoning → L3 Agent → L4 Innovation → L5 Organization.

\begin{importantbox}{Los niveles no guardan una relación lineal entre sí}
La intuición clave de Yang Zhilin (杨植麟) sobre L1-L5:
\begin{enumerate}
    \item Estos cinco niveles no son del todo secuenciales: algunas capacidades se desarrollan en paralelo
    \item El límite superior de Agent depende de la capacidad de Reasoning, pero Agent no tiene que esperar a que Reasoning madure por completo para empezar a desarrollarse
    \item El indicador de Innovation es que «el modelo participa en el desarrollo del modelo»; por ejemplo, K2 participa en el desarrollo de K3
    \item Ya ha aparecido un embrión de Organization: en los sistemas Multi-Agent, distintos Agent asumen distintos roles
\end{enumerate}
\end{importantbox}

\begin{knowledgebox}{Resolver los problemas de L3 con la tecnología de L4}
Hay una dependencia circular interesante: para resolver el problema de generalización de Agent (L3), puede hacer falta la tecnología de Innovation (L4), es decir, usar IA para entrenar y alinear IA. Esto muestra que entre estos niveles existe una relación de interacción compleja, y no una simple dependencia secuencial.
\end{knowledgebox}

\subsection{Resumen de la sección}

L1-L5 son varios hitos técnicos importantes, pero no guardan una relación estrictamente lineal. Reasoning y Agent son requisitos previos de Innovation y Organization, pero entre ellos existe una interdependencia compleja. La capacidad de cada nivel sigue mejorando de manera continua, sin un «tope» definido.
\section{De K1.5 a K2: la evolución técnica de Moonshot AI (月之暗面)}

\subsection{K1.5: validación de la técnica de aprendizaje por refuerzo}

K1.5 fue principalmente la validación de la ruta técnica de aprendizaje por refuerzo: Moonshot AI fue uno de los equipos que invirtieron en esta ruta desde etapas tempranas. Hallazgo clave:

\begin{importantbox}{El Reward de extremo a extremo supera al Process Reward}
Hallazgo técnico importante de K1.5: no es muy necesario un Process Reward Model (PRM) o una Value Function. Usar directamente un Reward de extremo a extremo permite lograr un entrenamiento muy bueno. El Process Reward incluso puede tener efectos secundarios durante el entrenamiento. Al principio, esta no era una conclusión muy clara.
\end{importantbox}

\subsection{K2: en busca de un mejor modelo base}

K2 tiene dos direcciones centrales en sus objetivos de diseño:

\subsubsection{Dirección uno: Token Efficiency, maximizar el valor de cada porción de datos}

\begin{importantbox}{Por qué Token Efficiency es más importante que Compute Efficiency}
El crecimiento de los datos de alta calidad es muy lento: puede aproximarse como una constante (por ejemplo, 30T Tokens de alta calidad). Bajo esta restricción:
\begin{itemize}
    \item \textbf{Compute Efficiency} (entrenar más rápido): aunque tiene valor, no eleva el límite superior de la inteligencia: solo se termina de entrenar más rápido con los mismos datos
    \item \textbf{Token Efficiency} (aprender mejor): con los mismos datos, el modelo absorbe más, la tasa de compresión es mayor y el loss baja más rápido: eleva directamente el límite superior de la inteligencia
\end{itemize}
\end{importantbox}

Moonshot AI mejora el Token Efficiency de tres maneras:

\begin{enumerate}
    \item \textbf{Optimizador Muon}: un nuevo optimizador basado en la ortogonalización de matrices, que sustituye a Adam, usado durante 10 años. En condiciones de Compute Optimal, el Token Efficiency mejora aproximadamente 2 veces: equivale a que aprender con 1 porción de datos sea comparable a que Adam aprenda con 2 porciones
    \item \textbf{Mayor dispersión}: al incorporar más parámetros mediante una mayor dispersión, cuantos más parámetros haya, mejor será la absorción de los mismos datos
    \item \textbf{Rephrase de datos}: reescribir los datos de alta calidad para que la misma porción de datos se aprenda varias veces en formas distintas, lo cual mejora la generalización
\end{enumerate}

\begin{knowledgebox}{Antecedentes técnicos del optimizador Muon}
El optimizador Muon fue propuesto por Keller. Su idea central es no considerar cada elemento de parámetro de la matriz de forma independiente (como hace Adam), sino tomar en cuenta la dependencia (dependency) entre los parámetros, y lograr un aprendizaje más eficiente mediante la ortogonalización de matrices.

La contribución de Moonshot AI fue llevar Muon a la práctica en modelos de lenguaje a gran escala. En una etapa temprana, el proyecto Moonlight lo validó por primera vez en un modelo de lenguaje de cierta escala; más adelante, al hacer Scale, surgieron nuevos problemas (como la explosión de MaxLogit), que se resolvieron proponiendo una nueva forma de Clipping.
\end{knowledgebox}

\begin{warningbox}{Reescribir datos no es lo mismo que copiarlos simplemente}
Aprender simplemente repitiendo la misma porción de datos provoca sobreajuste y problemas de generalización. La clave de la reescritura es introducir «nueva entropía»: hacer que los datos reescritos conserven el conocimiento central y, al mismo tiempo, tengan cierta diversidad, lo cual mejora la capacidad de generalización. Pero la elección de la forma de reescritura es crucial; cómo introducir nueva entropía de información sigue siendo una línea de investigación activa.
\end{warningbox}

\subsubsection{Dirección dos: capacidad Agentic y generalización}

\begin{importantbox}{La generalización de los Agent es el mayor desafío actual}
Para los modelos Agentic, el mayor desafío actual es la \textbf{generalización}:
\begin{itemize}
    \item Limitación de las técnicas de RL actuales: las tareas de entrenamiento y las métricas de evaluación suelen ser puntuales (por ejemplo, entrenar solo con SWE-Bench)
    \item Una mejora en la métrica no significa una mejora en la generalización: que la puntuación de SWE-Bench suba no implica que el modelo también sea mejor en otras tareas de código
    \item El problema de generalización de los Agent es más grave que en los modelos de conversación
    \item Los Benchmark son insuficientes o ya no son válidos
\end{itemize}
\end{importantbox}

Yang Zhilin (杨植麟) considera que una posible solución es «usar IA para entrenar IA»: hacer que el modelo participe en su propio proceso de entrenamiento y alineación, para mejorar la generalización de una manera más AI-Native.

\subsection{El proceso de investigación y desarrollo de K2}

\begin{itemize}
    \item La acumulación técnica comenzó un año antes (tecnologías como el optimizador Muon requieren experimentos de Scaling de ciclo largo para su validación)
    \item El ciclo de entrenamiento del propio K2 no fue largo, pero la preparación técnica previa requirió mucho tiempo
    \item El principal desafío encontrado: el optimizador Muon presentó el problema de explosión de MaxLogit durante el entrenamiento a gran escala (que no podía reproducirse en experimentos a pequeña escala)
    \item El equipo de investigación y el equipo de entrenamiento son el mismo equipo, porque los problemas durante el entrenamiento requieren una comprensión profunda de la tecnología subyacente para resolverse
\end{itemize}

\subsection{Resumen de la sección}

K1.5 validó la ruta de aprendizaje por refuerzo con Reward de extremo a extremo. K2 se centra en dos direcciones: (1) Token Efficiency (optimizador Muon + dispersión + reescritura de datos), (2) capacidad de generalización Agentic. La generalización de los Agent es el mayor cuello de botella actual, y quizá se necesite usar IA para entrenar IA para superarlo.

\section{Definición y futuro del Agent}

\subsection{Las dos características centrales del Agent}

\begin{importantbox}{Agent = multiturno + herramientas}
La definición que da Yang Zhilin (杨植麟) del Agent es concisa y esencial:
\begin{enumerate}
    \item \textbf{Multiturno} (Multi-turn): puede sostener interacciones múltiples; es una forma de Test-Time Scaling que permite al modelo completar tareas complejas
    \item \textbf{Herramientas} (Tool Use): la manera de conectar el cerebro con el mundo exterior; el motor de búsqueda conecta con internet y la capacidad de programar conecta con toda la automatización del mundo digital
\end{enumerate}
\end{importantbox}

\subsection{La generalidad es la capacidad que más le falta al Agent}

\begin{importantbox}{Del Agent general a las aplicaciones verticales}
Si el Agent tuviera una capacidad de generalización suficientemente buena, no harían falta tantos Agents verticales:
\begin{itemize}
    \item Un Agent general solo necesitaría conectarse a distintas herramientas (bases de datos, API o interfaces documentales a la medida) para completar tareas de un dominio vertical
    \item Pero justamente esa capacidad de generalización es lo que más le falta al Agent actual: si puede generalizar a herramientas y entornos que no vio durante el entrenamiento
    \item El sobreajuste a los Benchmark es un riesgo: subir la puntuación en unos cuantos Benchmark no representa la capacidad de generalización en escenarios reales
\end{itemize}
\end{importantbox}

\begin{knowledgebox}{El propósito del Agent no es simular a las personas}
Yang Zhilin (杨植麟) señala una distinción importante: el propósito del Agent es la generalidad (General Purpose), no simular a las personas. Los seres humanos también son generales (Universal Constructor); que la forma de actuar del Agent se parezca a la humana es solo una coincidencia, del mismo modo que el avión se diseñó como medio de transporte, no para volar como un ave.
\end{knowledgebox}

\subsection{Resumen de la sección}

Las dos características centrales del Agent son la interacción multiturno y el uso de herramientas. El mayor reto actual es la capacidad de generalización: si puede desempeñarse bien en herramientas y entornos que no ha visto. El objetivo del Agent es la inteligencia general, no simular la conducta humana.

\section{Scaling Law y arquitectura de modelos}

\subsection{El nuevo cuello de botella del Scaling de preentrenamiento}

\begin{warningbox}{Los datos de alta calidad se acercan a una constante}
La velocidad de crecimiento de los datos de entrenamiento de alta calidad es mucho menor que la del crecimiento del cómputo:
\begin{itemize}
    \item El volumen de datos de texto de alta calidad puede considerarse aproximadamente constante
    \item Los datos multimodales no logran elevar de forma efectiva el «coeficiente intelectual» del texto
    \item Esto hace que la Token Efficiency se vuelva una dirección de optimización más importante que la Compute Efficiency
\end{itemize}
\end{warningbox}

\subsection{El optimizador Muon: rompiendo 10 años de dominio de Adam}

El optimizador Adam dominó el entrenamiento de aprendizaje profundo durante casi 10 años. El optimizador Muon, al considerar las relaciones de dependencia estructural entre matrices de parámetros, logra una mejora de aproximadamente 2 veces en la Token Efficiency:

\begin{itemize}
    \item Adam: considera cada elemento de parámetro de forma independiente, ignorando las relaciones estructurales entre parámetros
    \item Muon: considera la dependency entre parámetros matriciales, optimizando mediante ortogonalización
    \item Efecto: el efecto de aprendizaje de 30T tokens de alta calidad equivale al de Adam aprendiendo 60T tokens
\end{itemize}

\subsection{Hallazgos sobre el Scaling de RL}

\begin{knowledgebox}{Diseño de Curriculum en aprendizaje por refuerzo}
En el entrenamiento de agentes, no se puede dejar que el modelo resuelva de entrada la tarea más difícil (como demostrar una conjetura matemática no resuelta), pues de lo contrario la Sample Efficiency es extremadamente baja. Una mejor forma es:
\begin{itemize}
    \item Usar tareas de dificultad media para el entrenamiento
    \item Elevar la dificultad de forma gradual (un proceso de escalar una montaña)
    \item Combinarlo con una buena estrategia de Sampling
\end{itemize}
En esencia, esto es un proceso de Curriculum Learning.
\end{knowledgebox}

\subsection{Resumen de la sección}

El nuevo cuello de botella del Scaling de preentrenamiento es la disponibilidad limitada de datos de alta calidad. El optimizador Muon, mediante la ortogonalización de matrices, rompe los 10 años de dominio de Adam y logra una mejora de 2 veces en la Token Efficiency. El entrenamiento de RL requiere un diseño de Curriculum para garantizar la Sample Efficiency.


\section{El dilema de la evaluación en el entrenamiento de modelos}

\subsection{Fallo de los Benchmarks y generalización del Agent}

\begin{warningbox}{Riesgo de sobreajuste en los Benchmarks}
El campo de los Agent enfrenta actualmente un problema serio de evaluación:
\begin{itemize}
    \item La cantidad de Benchmarks de Agent disponibles es limitada
    \item Elevar la puntuación de un Benchmark no representa la capacidad real: muchas veces se trata de un sobreajuste a una tarea específica
    \item La percepción de los usuarios en escenarios OOD (Out-of-Distribution) se desliga de la puntuación del Benchmark
    \item «Quien siembra melones cosecha melones, quien siembra frijoles cosecha frijoles»: se mejora aquello para lo que se entrena, pero la generalización es limitada
\end{itemize}
\end{warningbox}

\subsection{El origen del dilema de evaluación}

Yang Zhilin (杨植麟) considera que el origen del dilema de evaluación es el siguiente:

\begin{enumerate}
    \item El aprendizaje por refuerzo depende en última instancia de una buena evaluación o verificación: los problemas matemáticos y la programación con casos de prueba (Test Case) pueden verificarse de manera automática, pero para las tareas complejas de extremo a extremo es difícil encontrar una buena forma de evaluación
    \item Los nuevos problemas que trae la mejora de la capacidad de los modelos: la complejidad de las tareas aumenta y la forma de evaluación necesita mantenerse al día
    \item La característica de distribución de cola larga de las tareas de Agent: es difícil cubrir todos los escenarios reales
\end{enumerate}

\subsection{Resumen de la sección}

El fallo de los Benchmarks y la generalización insuficiente de los Agent son el dilema de evaluación más destacado en la actualidad. El origen está en que las tareas complejas son difíciles de verificar de manera automática, así como en el Gap entre la distribución de entrenamiento y los escenarios reales de uso. Hacer que la IA participe en la evaluación y el entrenamiento (la IA entrena a la IA) podría ser una dirección de avance.

\section{Producto propio vs. ecosistema de terceros}

\subsection{Dos paradigmas de producto}

\begin{importantbox}{La lógica de las empresas de modelos que hacen producto propio}
\begin{itemize}
    \item \textbf{Paradigma de terceros}: dado un modelo, se hace ingeniería inversa de la mejor forma de usarlo (System Prompt, diseño de herramientas, Context Engineering); en esencia, se busca ajustar la distribución de entrenamiento del modelo
    \item \textbf{Paradigma propio}: primero se diseñan las herramientas y el entorno, y después se entrena el modelo en ese entorno. El modelo rinde mejor de forma natural en su propio entorno, sin necesidad de ingeniería inversa
\end{itemize}
Cloud Code y ChatGPT Agent son representantes del producto propio. Moonshot AI (月之暗面) todavía se centra sobre todo en el modelo, pero considera que el producto propio será una tendencia importante.
\end{importantbox}

\subsection{Resumen de la sección}

El producto propio, mediante el diseño directo del entorno y el entrenamiento del modelo, tiene un techo más alto; el producto de terceros, mediante la ingeniería inversa de la forma de usar el modelo, es flexible pero está limitado por la distribución de entrenamiento del modelo. En el futuro, los dos paradigmas coexistirán de forma complementaria.

\section{De K2 a la siguiente generación de Agentes: cuáles problemas son los más difíciles}

\subsection{Token Efficiency y las restricciones del sistema}
En la entrevista hay un juicio muy explícito: la competencia entre los grandes modelos del futuro no será solo «quién tiene más compute», sino quién puede consumir tokens de alta calidad de manera más eficiente. Yang Zhilin (杨植麟) coloca en el mismo panorama el optimizador Muon, la reescritura de datos, los datos sintéticos y la capacidad de escalamiento en tiempo de prueba; en esencia, está respondiendo una pregunta: cuando los datos de entrenamiento de alta calidad son cada vez más escasos, ¿cómo lograr que cada token produzca un mayor rendimiento?

\begin{importantbox}{Por qué Token Efficiency es más crucial que la aparente Compute Efficiency}
\begin{itemize}
    \item los datos realmente de alto valor ya son escasos de por sí; no se puede esperar resolver el problema apilando datos sin límite;
    \item entrenar tareas de Agent tiene un costo más alto, porque además hay que pagar el costo de rollout, tool use, ejecución del entorno y evaluación;
    \item si la tasa de aprovechamiento de tokens no es alta, incluso con una capacidad de cómputo mayor, el sistema se detendrá frente al cuello de botella de los datos de alta calidad.
\end{itemize}
\end{importantbox}

\subsection{Los tres umbrales duros después de K2}
Si se condensa el juicio de la entrevista en una formulación más orientada a la ingeniería, después de K2 quedan al menos tres umbrales que aún no se han cruzado de verdad: el primero es la generalización de los Agentes, el segundo es el sistema de evaluación, y el tercero es el ciclo cerrado a nivel de sistema.

\begin{table}[H]
\centering
\begin{tabular}{p{2.8cm}p{4.2cm}p{4.5cm}}
\toprule
\textbf{Umbral} & \textbf{Juicio en la entrevista} & \textbf{Por qué es difícil} \\
\midrule
Generalización de los Agentes & Una puntuación alta en el Benchmark no significa que el escenario real se resuelva correctamente & La distribución de tareas tiene una cola larga; la distribución de entrenamiento no coincide con el uso real \\
Sistema de evaluación & Las tareas complejas de extremo a extremo carecen de una forma de verificación automática & Sin un verifier estable, es muy difícil hacer RL de alta calidad \\
Ciclo cerrado del sistema & Un producto propio integra con más facilidad el entorno y el entrenamiento & El ecosistema de terceros depende de la ingeniería inversa para ajustarse a la distribución de entrenamiento del modelo \\
\bottomrule
\end{tabular}
\caption{Lo verdaderamente difícil después de K2 no es «si se puede hacer una demostración», sino cómo escalar de manera estable hacia tareas reales}
\end{table}

\begin{warningbox}{Lo que más fácilmente se subestima no es la capacidad del modelo, sino el costo de la evaluación}
La crítica de la entrevista al Benchmark es muy contundente: las tareas de Agent no son como los problemas de matemáticas o las tareas de código con test case; mucho trabajo real no tiene una respuesta correcta natural. Sin una evaluación confiable, no se puede hacer aprendizaje por refuerzo de manera estable; sin aprendizaje por refuerzo, es muy difícil interiorizar en el modelo las capacidades de cadena larga. Esta es una razón importante por la que la velocidad de iteración de los Agentes es más lenta de lo que el mundo exterior imagina.
\end{warningbox}

\subsection{Resumen de la sección}
Si K2 representa un avance en la capacidad del modelo, lo que realmente determinará el techo en la siguiente etapa será la eficiencia de los tokens, la evaluación confiable y el ciclo cerrado sistemático del producto. Precisamente por eso, a lo largo de toda la entrevista Yang Zhilin coloca de manera constante «modelo», «datos», «evaluación» y «producto» dentro del mismo panorama de discusión.

\subsection{Implicaciones para las organizaciones y los equipos de producto}
Esta entrevista también ofrece enseñanzas directas más allá del equipo de investigación. Si en el futuro los productos propios serán cada vez más importantes, entonces la forma de la organización también cambiará: el equipo de modelos no puede encargarse solo del entrenamiento, y el equipo de producto tampoco puede encargarse solo de envolver una capa de UI; ambos necesitan entender el entorno de herramientas, el ciclo cerrado de evaluación y las tareas reales del usuario.

\begin{knowledgebox}{Por qué la ruta del producto propio obliga a reestructurar la organización}
\begin{itemize}
    \item hacer un producto propio implica que el diseño del entorno, el entrenamiento del modelo y la interacción del producto deben planearse de manera sincronizada;
    \item si la organización sigue completamente dividida en «grupo de modelos» y «grupo de producto», mucha retroalimentación clave no puede regresar al entrenamiento;
    \item cuanto más se adentra un producto de Agent en el flujo de trabajo, más necesita conectar la telemetry del producto, los casos de falla y los datos de entrenamiento.
\end{itemize}
\end{knowledgebox}

\section{Metodología de emprendimiento que se refleja en la entrevista}

\subsection{Por qué una empresa de modelos no puede dedicarse solo a los modelos}
Aunque toda la entrevista gira en torno a K2, Muon, RL y los Agent, lo que Yang Zhilin (杨植麟) reitera en realidad es un juicio más amplio: \textbf{la capacidad del modelo, el sistema de evaluación y el entorno del producto deben evolucionar de manera sincronizada}. Si se entiende la empresa solo como «entrenar un modelo base más fuerte», se subestima el tramo más difícil que viene después, porque lo que en verdad convierte la capacidad en valor para el usuario es el sistema que rodea al modelo.

\begin{importantbox}{Del suministro de capacidad al cierre de la tarea: la cadena de valor se está desplazando hacia atrás}
\begin{itemize}
    \item El preentrenamiento y el entrenamiento posterior determinan el límite superior de la capacidad del modelo;
    \item el sistema de evaluación determina si esas capacidades pueden identificarse de manera estable e iterarse de forma continua;
    \item el entorno del producto determina si el modelo puede convertir esa capacidad en resultados dentro de flujos de trabajo reales.
\end{itemize}
Por eso en la entrevista aparecen a la vez temas como K2, Agent, Benchmark y los productos propios de la empresa. No son puntos de interés dispersos, sino distintos eslabones de una misma cadena de valor.
\end{importantbox}

Yang Zhilin no entiende el «producto» como una capa de interfaz de chat que se agrega al final, sino algo más cercano a «construirle al modelo un entorno en el que pueda aprender y ejecutar tareas a largo plazo». Bajo esta perspectiva, el producto es a la vez canal de distribución y parte del entorno de entrenamiento. Si una empresa de modelos pasa esto por alto, es fácil que se vea muy fuerte a nivel de demo, pero que no logre construir un interés compuesto en tareas reales.

\subsection{El ritmo y la ventana de oportunidad de las empresas emergentes}
Esta entrevista también deja entrever un juicio sobre el ritmo de una empresa emergente: en un periodo de cambio de paradigma tan intenso, el equipo debe priorizar las capacidades básicas capaces de generar un interés compuesto, en lugar de perseguir tendencias de corto plazo. Ya sea el optimizador Muon, la eficiencia de tokens, la generalización de los Agent o el sistema de evaluación, en el fondo todos compiten por la pendiente tecnológica de los próximos años.

\begin{knowledgebox}{Tres capas de ventana de oportunidad que se pueden abstraer de la entrevista}
\begin{enumerate}
    \item \textbf{Ventana algorítmica}: quien encuentre primero una mayor eficiencia de tokens, un RL más estable y mejores verifiers podrá ampliar su capacidad a un costo menor;
    \item \textbf{Ventana de sistema}: quien logre conectar entre sí las herramientas, el entorno, la memoria y la evaluación tendrá más facilidad para convertir la capacidad del modelo en experiencia de producto;
    \item \textbf{Ventana organizacional}: quien logre que la investigación, la ingeniería y los datos del producto formen un ciclo de retroalimentación rápido tendrá más facilidad para iterar de manera continua.
\end{enumerate}
\end{knowledgebox}

\begin{warningbox}{El error de juicio más peligroso es confundir una ventaja de corto plazo con una barrera de largo plazo}
La cautela que muestra la entrevista frente a los Benchmark también aplica al juicio sobre una empresa. Un efecto de demostración de corto plazo, la mejora en un benchmark puntual, o incluso el éxito repentino de un producto, no constituyen necesariamente una barrera de largo plazo. La verdadera barrera proviene de: poder absorber de manera continua datos de alta calidad, poder construir una evaluación confiable y poder retroalimentar las tareas del usuario hacia el ciclo cerrado de entrenamiento y producto.
\end{warningbox}

\textit{«Los problemas son inevitables, y los problemas también se pueden resolver.»} Si se traslada esta frase al contexto del emprendimiento, significa que el equipo no debe esperar que un «modelo definitivo» resuelva todos los problemas de una sola vez, sino que debe construir, en torno a los cuellos de botella clave, un mecanismo de ascenso continuo. K2 es solo un punto alto en esta etapa, no el final del camino.

\subsection{Resumen de la sección}
La metodología de emprendimiento que se puede abstraer de esta entrevista es: no hay que ver por separado el modelo, la evaluación y el producto, sino considerarlos como distintas capas de un mismo sistema. El equipo que gane la competencia de largo plazo en el futuro no será solo el que entrene un modelo fuerte, sino el que logre convertir de manera continua ese modelo fuerte en el cierre de tareas reales.
\section{Síntesis y ampliación}

\subsection{Resumen de las ideas centrales}

\begin{importantbox}{Puntos clave de la entrevista de Yang Zhilin (杨植麟)}
\begin{enumerate}
    \item \textbf{La IA es una montaña de nieve sin fin}: los problemas son inevitables, pero se pueden resolver. Cada problema resuelto equivale a ascender unos cientos de metros, y al mismo tiempo aparecen problemas nuevos
    \item \textbf{Dos tipos de Test-Time Scaling}: los modelos de razonamiento (cerebro en una cubeta, hacen Scale de los tokens de pensamiento) y los agentes (interactúan con el mundo, hacen Scale de los turnos)
    \item \textbf{Token Efficiency es más importante que Compute Efficiency}: los datos de alta calidad son el cuello de botella, y el optimizador Muon logra una mejora de 2 veces
    \item \textbf{La generalización de los agentes es el mayor reto}: hay un sobreajuste severo a los Benchmark, y tal vez se necesite que una IA entrene a otra IA para superarlo
    \item \textbf{AGI es una dirección, no un punto final}: no habrá un momento en que de pronto se alcance la AGI, sino una mejora continua de las distintas capacidades
    \item \textbf{Tendencia de los productos de primera parte}: las empresas de modelos avanzan hacia diseñar el entorno y entrenar el modelo a la vez, lo que eleva el límite superior
\end{enumerate}
\end{importantbox}

\subsection{Reflexiones adicionales}

\begin{itemize}
    \item \textbf{La automejora de la IA}: cuando K2 pueda participar en el desarrollo de K3, ¿qué tipo de ciclo de aceleración se formará?
    \item \textbf{Renovación del paradigma de evaluación}: ¿cómo diseñar métodos de evaluación de agentes que no se sobreajusten?
    \item \textbf{Superar el cuello de botella de datos}: además de reescribir y sintetizar, ¿qué otras vías existen para mejorar el Token Efficiency?
    \item \textbf{El umbral de los agentes generales}: ¿qué tan fuerte debe ser la capacidad de generalización de un agente para que realmente sustituya las soluciones verticales?
\end{itemize}

\subsection{Lecturas adicionales}

\begin{itemize}
    \item Informe técnico de K2 (publicado oficialmente por Moonshot AI (月之暗面))
    \item David Deutsch, \textit{The Beginning of Infinity}
    \item Moonlight: el optimizador Muon en modelos de lenguaje a gran escala
    \item El sistema de niveles L1--L5 de OpenAI
    \item El diseño de productos de primera parte de Claude Code y ChatGPT Agent
\end{itemize}

%% --- Fin del contenido --- %%

\end{document}
